\subsection{Mixed-Integer Convex Programming Formulation via Perspective Relaxation}

Directly solving the problem as a Mixed-Integer Nonlinear Programming (MINLP) formulation poses significant computational challenges due to the presence of non-convex constraints, particularly the bilinear relationship between the binary decision variables $y_e$ and the continuous state variables $x_u$. To overcome this and transform the problem into a tractable Mixed-Integer Convex Programming (MICP) form, we apply the perspective relaxation technique. The primary motivation of this method is to isolate the non-convexity by ``lifting'' the continuous variables into a higher-dimensional space. Instead of utilizing global variables $x_u$, we introduce local proxy variables $z_{e,u}, z_{e,v} \in \mathbb{R}^n$ for each edge $e=(u,v)$, representing the state contribution of the corresponding vertices when that edge is active.

For each edge $e=(u, v)$, the local proxy variables $z_{e,u}$ and $z_{e,v}$ specifically model the contribution of the state at vertices $u$ and $v$ to the traversal of edge $e$. If the edge is not active (i.e., $y_e=0$), these contributions are forced to zero. This relationship is expressed as:
\[
z_{e,u} \in y_e X_u \quad \text{and} \quad z_{e,v} \in y_e X_v
\]
which implies:
\[\begin{cases}
z_{e,u} \in X_u & \text{if } y_e = 1 \\
z_{e,u} = 0 & \text{if } y_e = 0
\end{cases}
\]

This logic is formalized using the perspective function $\tilde{\ell}_e((z_{e,u}, z_{e,v}), y_e)$, which is the perspective of the original cost function $\ell_e$. A fundamental property in convex analysis states that if $\ell_e$ is convex, its perspective $\tilde{\ell}_e$ is jointly convex in all arguments. This property is crucial as it enables the use of efficient convex solvers. The perspective function is defined as:
\[
\tilde{\ell}_e(z_{e,u}, z_{e,v}, y_e) = 
\begin{cases} 
y_e \ell_e\left(\frac{z_{e,u}}{y_e}, \frac{z_{e,v}}{y_e}\right) & \text{if } y_e > 0 \\
0 & \text{if } y_e = 0 \text{ and } z_{e,u} = z_{e,v} = 0 \\
+\infty & \text{otherwise}
\end{cases}
\]

Crucially, rather than enforcing explicit equality constraints such as $z_{e,u} = y_e x_u$, which are non-convex, we extend the standard flow conservation law~\eqref{eq:flow_conservation} to the continuous variables. This mandates that for any node $u$ (excluding the source and target), the sum of continuous states entering the node via incoming edges must equal the sum of continuous states leaving the node via outgoing edges.

The resulting MICP formulation for the Shortest Path Problem in GCS is defined as follows:

\begin{subequations}
\begin{align}
\text{minimize} \quad & \sum_{e=(u,v) \in \mathcal{E}} \tilde{\ell}_e((z_{e,u}, z_{e,v}), y_e) \label{eq:gcs_obj_perspective} \\
\text{subject to} \quad & \sum_{e \in \mathcal{E}^+(u)} y_e - \sum_{e \in \mathcal{E}^-(u)} y_e = \delta_u, \quad \forall u \in \mathcal{V} \label{eq:binary_flow} \\
& \sum_{e \in \mathcal{E}^+(u)} z_{e,u} - \sum_{e \in \mathcal{E}^-(u)} z_{e,v} = 
\begin{cases}
x_s & \text{if } u = s \\
-x_t & \text{if } u = t \\
0 & \text{otherwise}
\end{cases}, \quad \forall u \in \mathcal{V} \label{eq:spatial_flow} \\
& (z_{e,u}, z_{e,v}) \in y_e \mathcal{X}_e, \quad \forall e \in \mathcal{E} \label{eq:perspective_edge} \\
& z_{e,u} \in y_e X_u, \quad z_{e,v} \in y_e X_v, \quad \forall e=(u,v) \in \mathcal{E} \label{eq:perspective_vertex} \\
& y_e \in \{0, 1\}, \quad \forall e \in \mathcal{E} \label{eq:binary_constraint}
\end{align}
\end{subequations}

In this model, Equation~\eqref{eq:binary_flow} ensures topological connectivity, where $\delta_u$ is defined as $1$ for the source, $-1$ for the target, and $0$ otherwise. Equation~\eqref{eq:spatial_flow} ensures that the continuous trajectory is seamless across the graph structure. Equations~\eqref{eq:perspective_edge} and~\eqref{eq:perspective_vertex} enforce geometric containment within scaled convex sets; specifically, if $y_e=0$, the associated continuous variables are constrained to the origin, incurring zero cost. This formulation provides a computationally efficient and tight relaxation, allowing the solver to identify near-global optimal trajectories that satisfy both the discrete graph topology and the complex continuous constraints inherent in motion planning.

\subsection{Convexification of the GCS Problem via the RLT Technique}

To address the inherent non-convexity of the initial Mixed-Integer Non-Convex Program (MINCP)—specifically the bilinear constraints coupling binary decision variables with continuous state variables—we employ a convexification strategy based on the Reformulation-Linearization Technique (RLT) of the first level. The fundamental objective of this approach is to replace intractable constraints with tight convex envelopes. This ensures that the resulting Mixed-Integer Convex Program (MICP) preserves the integer feasible set of the original problem while providing strong lower bounds when the integrality constraints are relaxed.

\subsubsection{Theoretical Foundation: The Duality between Perspective Cones and Valid Inequalities}

Before detailing the optimization model, we must establish the mathematical connection between perspective cones and linear inequalities. Consider the dual cone of the perspective cone $\tilde{X}$, denoted as $(\tilde{X})^*$. This is defined as the set of vectors $(a, b)$ in the extended space such that their inner product with every element in the cone $\tilde{X}$ is non-negative:
\begin{equation}
    (\tilde{X})^* := \left\{ (a, b) \in \mathbb{R}^{n+1} \mid (a, b)^\top k \ge 0, \quad \forall k \in \tilde{X} \right\}.
\end{equation}

A critical result utilized in this formulation is the equivalence between the dual cone $(\tilde{X})^*$ and the set of valid linear inequalities for the set $X$. Geometrically, this represents the intersection of all closed half-spaces containing $X$. We formalize this in the following proposition:

\textbf{Proposition.} The dual cone $(\tilde{X})^*$ is precisely the set of coefficients $(a, b)$ that define valid linear inequalities over the set $X$; that is, $(\tilde{X})^* \equiv X^\circ$.

\textit{Proof.}
To demonstrate this, consider an arbitrary element $k \in \tilde{X}$. By the definition of a perspective cone, $k$ can be represented as $k = (\lambda x, \lambda)$ for some $x \in X$ and a scalar $\lambda \ge 0$. The condition for a vector pair $(a, b)$ to belong to the dual cone $(\tilde{X})^*$ can be expanded as follows:
\begin{align}
    (a, b) \in (\tilde{X})^* &\iff (a, b)^\top (\lambda x, \lambda) \ge 0, \quad \forall x \in X, \forall \lambda \ge 0 \nonumber \\
    &\iff \lambda (a^\top x + b) \ge 0, \quad \forall x \in X, \forall \lambda \ge 0.
\end{align}
Observe that the final inequality must hold for all non-negative values of $\lambda$. In the case where $\lambda > 0$, we may divide both sides by $\lambda$ without altering the inequality sign, yielding the condition $a^\top x + b \ge 0$. The case $\lambda = 0$ is trivially satisfied. Consequently, the necessary and sufficient condition for $(a, b) \in (\tilde{X})^*$ is that the linear expression $a^\top x + b$ remains non-negative over the entire set $X$. This confirms that $(\tilde{X})^*$ constitutes the set of coefficients for all valid inequalities on $X$, completing the proof. \qed

\subsubsection{Construction of Convex Constraints}

Leveraging the proposition above, we can ``lift'' any valid linear constraint from the flow variable space into the mixed state-flow space. Consider an arbitrary vertex $v$ and a valid linear inequality concerning the flow variables $y_e$ incident to $v$, given by:
\begin{equation}
    \sum_{e \in E_v} c_e y_e + d \ge 0, \label{eq:valid_ineq}
\end{equation}
where $E_v$ denotes the set of edges adjacent to $v$, and $d$ is a problem-dependent constant. By applying the perspective principle, we transform the inequality (~\ref{eq:valid_ineq}) into a constraint within the convex cone $\tilde{X}_v$. To facilitate this, we introduce proxy variables $z_e$ and $z'_e$. The variable $z_e$ represents the state contribution at the source vertex of edge $e$ when active, while $z'_e$ represents the state contribution at the target vertex.

This transformation yields the following convex cone constraint:
\begin{equation}
    \left( \sum_{e \in E_{\text{in}}(v)} c_e z'_e + \sum_{e \in E_{\text{out}}(v)} c_e z_e + d x_v, \quad \sum_{e \in E_v} c_e y_e + d \right) \in \tilde{X}_v. \label{eq:lifted_constraint}
\end{equation}

Equation (~\ref{eq:lifted_constraint}) constrains the relationship between the intermediate state variables ($z, z'$) and the flow variables $y$, ensuring they reside within the perspective cone of the feasible set $X_v$.

A specific yet crucial application of this Lemma involves the non-negativity constraint of the flow, i.e., $y_e \ge 0$. For an edge $e=(u, v)$:
\begin{itemize}
    \item At the source vertex $u$, applying the lifting with $c_e=1$ yields the constraint $(z_e, y_e) \in \tilde{X}_u$.
    \item At the target vertex $v$, applying the lifting with $c_e=1$ yields the constraint $(z'_e, y_e) \in \tilde{X}_v$.
\end{itemize}
These constraints effectively function as an ``on/off'' mechanism: if $y_e=1$, the proxy variable $z_e$ is constrained to the set $X_u$; conversely, if $y_e=0$, $z_e$ is forced to the origin.

\subsubsection{Complete MICP Formulation}

Synthesizing the preceding analysis, we propose the Mixed-Integer Convex Programming (MICP) formulation for the Shortest Path Problem on GCS, as described in Theorem 5.7. This model completely eliminates bilinear components, replacing them with convex cone constraints:

\begin{subequations}~\label{eq:micp_formulation}
\begin{align}
  \underset{y, z, z'}{\text{minimize}} \quad & \sum_{e \in E} \tilde{\ell}_e((z_e, z'_e), y_e) \tag{5.5a} \\
  \text{subject to} \quad & \sum_{e \in E_{\text{out}}(s)} y_e = 1, \quad \sum_{e \in E_{\text{in}}(t)} y_e = 1, \tag{5.5b} \\
  & \sum_{e \in E_{\text{out}}(v)} y_e \le 1, \quad \forall v \in V \setminus \{s, t\}, \tag{5.5c} \\
  & \sum_{e \in E_{\text{in}}(v)} (z'_e, y_e) = \sum_{e \in E_{\text{out}}(v)} (z_e, y_e), \quad \forall v \in V \setminus \{s, t\}, \tag{5.5d} \\
  & (z_e, y_e) \in \tilde{X}_u, \quad (z'_e, y_e) \in \tilde{X}_v, \quad \forall e = (u, v) \in E, \tag{5.5e} \\
  & y_e \in \{0, 1\}, \quad \forall e \in E. \tag{5.5f}
\end{align}
\end{subequations}

The formulation (\ref{eq:micp_formulation}) is composed of four primary components that ensure graph connectivity and geometric consistency.

First, the objective function (5.5a) represents the total cost across all edges, where $\tilde{\ell}_e$ is the perspective function of the original cost metric. Since perspective operations preserve convexity, the optimization problem maintains a tractable convex structure.

Second, constraints (5.5b) and (5.5c) enforce classical flow conservation, guaranteeing the existence of a path from the source $s$ to the target $t$ while restricting each intermediate vertex to be visited at most once.

Third, constraint (5.5d) represents the spatial flow conservation, a distinct feature compared to standard shortest path problems. This constraint is derived by multiplying the flow conservation equation at vertex $v$ by the state variable $x_v$. Physically, it ensures trajectory continuity: the aggregated state ``entering'' vertex $v$ (via $z'_e$) must equal the aggregated state ``leaving'' vertex $v$ (via $z_e$).

Finally, constraints (5.5e) result from the convexification process discussed previously. They serve as the convex envelopes for the proxy variables $z_e$ and $z'_e$, ensuring that when an edge $e=(u,v)$ is selected ($y_e=1$), these variables correspond to valid points within the sets $X_u$ and $X_v$, respectively.